\chapter{Background and Motivation}
\label{ch:background}

The discourse surrounding AI ethics is not new, but its urgency has intensified as AI systems have transitioned from theoretical research to impactful, real-world applications. The motivation for this project stems from three interconnected observations: the escalating scale of AI-driven societal harm, the inadequacy of purely principles-based governance, and the complex, multi-stakeholder nature of ethical decision-making.

First, the potential for AI to cause significant, often unintended, harm is now well-documented. High-profile examples include biased facial recognition systems that exhibit higher error rates for women and people of color~\cite{buolamwini2018gender}, automated hiring tools that perpetuate historical gender and age biases~\cite{eeoc2022itutorgroup}, and opaque credit scoring algorithms that deny loans without clear justification~\cite{oneil2016weapons}. These failures are not merely technical; they have profound social, economic, and human rights implications. They underscore the reality that technical performance metrics, such as accuracy, are insufficient for evaluating the overall acceptability of an AI system. A broader, more holistic evaluative lens is required, one that explicitly accounts for societal values like fairness, accountability, and transparency.

Second, while the proliferation of ethical AI principles from governments and corporations is a positive development, these high-level declarations often fail to provide actionable guidance for practitioners~\cite{hagendorff2020ethics}. A principle such as \enquote{AI should be fair} is a necessary starting point, but it does not specify what fairness means in a given context, how it should be measured, or how it should be balanced against other competing values like accuracy or security. This \enquote{principles-to-practice} gap leaves developers without the tools to navigate complex ethical trade-offs during the design and implementation process~\cite{mittelstadt2019principles}. The result is often an ad-hoc, inconsistent, and unauditable approach to ethics, where critical decisions are made without a formal, defensible methodology.

Third, ethical questions in AI are rarely, if ever, one-dimensional. Different stakeholders, including developers, business owners, regulators, and the communities affected by the AI system, hold different values, priorities, and tolerances for risk~\cite{freeman1984stakeholder}. A developer might prioritize system robustness and security, a regulator might focus on legal compliance and accountability, and an affected community member might be most concerned with fairness and the right to redress. A decision that appears optimal from one perspective may be unacceptable from another. This inherent multi-stakeholder conflict is a central challenge in AI governance~\cite{mitchell1997stakeholder}. Ignoring it leads to solutions that fail to achieve broad legitimacy and may ultimately be rejected by society. Therefore, a functional ethical framework must not only evaluate AI systems but also provide a mechanism for surfacing, understanding, and negotiating these diverse stakeholder perspectives.

This project is motivated by the conviction that a computational approach can help bridge these gaps. By creating a structured, data-driven decision framework, MEDF aims to make ethical evaluation more rigorous, transparent, and inclusive. It seeks to empower AI development teams with a practical tool that translates abstract principles into concrete analysis, facilitates a more democratic dialogue about ethical trade-offs, and ultimately contributes to the creation of more responsible and trustworthy AI systems.
